\documentclass{beamer}

\usetheme{amber}

\usepackage{polyglossia}
\setmainlanguage{russian}
\setmainfont{Liberation Serif}
\setsansfont{Fira Sans}
\setmonofont{Fira Code}

\usepackage{mathpartir}

\title{Системный язык программирования с зависимыми типами}
\author{Соколов Павел}
\date{14 декабря 2022}

\begin{document}

\frame{\titlepage}

\begin{frame}{Мотивация}
  Все снова хотят, чтобы было быстро, но при этом хотят, чтобы было безопасно.
  \begin{itemize}
    \item C --- не позволяет выстраивать правильный уровень абстракции;
    \item C++ --- слишком большой и небезопасный язык;
    \item Go --- спайки из-за уборки мусора; не позволяет выстраивать правильный
      уровень абстракции.
    \item Rust --- сложные паттерны управления памятью невыразимы без unsafe и
      без просадки в производительности.
  \end{itemize}
  Мы можем лучше.
\end{frame}

\begin{frame}{Чуть менее известные решения}
  \begin{itemize}
    \item Flux --- refinement types для Rust. Громоздкие аннотации, не решает
      фундаментальную проблему (пока?).
    \item Nim --- инновационный сборщик циклов, ORC. Всё ещё заметен оверхед.
    \item Koka --- Functional-but-in-place (FBIP). Не системный язык
      программирования.
  \end{itemize}
\end{frame}

\section{Зависимые типы}

\frame{\sectionpage}

\begin{frame}{Определение}
  В диалектах лямбда-исчисления могут быть следующие зависимости:

  \begin{itemize}
    \item Термов от термов --- $\lambda x. M(x)$;
    \item Термов от типов --- $\Lambda A. M(A)$;
    \item Типов от типов --- $\lambda A. T(A)$;
    \item Типов от термов --- $\Pi x. T(x)$.
  \end{itemize}

  Формально, только последняя зависимость называется зависимыми типами, но в
  таких системах, как правило, есть и все остальные зависимости (исключение ---
  Twelf).
\end{frame}

\begin{frame}{Проверка типов}
  Чтобы существовала программная реализация системы, трёхместное отношение
  типизации $\Gamma \vdash M : \tau$ должно быть разрешимым.

  В частности, необходимо программным образом проверять следующее правило:

  \begin{mathpar}
    \inferrule
      {\Gamma \vdash f : \Pi (x : A) . B \\ \Gamma \vdash t : A}
      {\Gamma \vdash f \; t : B[x := t]}
  \end{mathpar}

  Угадать правильный $x$ для рекурсивного вызова от $f$ не получится, так что
  наивный алгоритм проверки типов не подойдёт, необходимо задействовать вывод
  типов.
\end{frame}

\begin{frame}{Вывод типов}
  Задача вывода типов звучит так: пусть дан контекст $\Gamma$ и терм $t$.
  Вывести все типы $\tau$ такие, что $\Gamma \vdash t : \tau$.

  Вернёмся к правилу с предыдущего слайда:

  \begin{mathpar}
    \inferrule
      {\Gamma \vdash f : \Pi (x : A) . B \\ \Gamma \vdash t : A}
      {\Gamma \vdash f \; t : B[x := t]}
  \end{mathpar}

  Пусть мы вывели тип $f$ и тип $t$. Осталось проверить, что вхождения $A$ в оба
  типа одинаковы. То есть, необходимо проверить равенство типов.

  \begin{block}{Двусторонний вывод}
    Либо можно использовать двусторонний вывод типов, но там, опять же,
    необходимо проверять равенство типов:

    \begin{mathpar}
      \inferrule
        {\Gamma \vdash t \Rightarrow A \\ A = B}
        {\Gamma \vdash t \Leftarrow B}
    \end{mathpar}
  \end{block}
\end{frame}

\begin{frame}[fragile=singleslide]{Равенство типов}
  \begin{block}{Задача}
    В более простых системах равенство типов является структурным. Однако в
    зависимых типах всё не так просто: пусть \verb|Vect a n| --- вектор из $n$
    значений типа $a$. Одинаковы ли типы \verb|Vect a (n + m)| и
    \verb|Vect a (m + n)|?
  \end{block}
  \begin{block}{Решение}
    Полного решения нет. Варианты:
    \begin{itemize}
      \item Ограничить термы, от которых можно зависеть, термами,
        вычислимыми на этапе компиляции*;
      \item Ограничить термы, от которых можно зависеть, термами,
        первопорядковые теории равенства которых разрешимы;
      \item Разрешить пользователю приравнивать доказуемо равные термы.
    \end{itemize}
  \end{block}
\end{frame}

\begin{frame}{Системы типов Мартина-Лёфа}
  В последнем случае необходимо ввести т.н. ``тип равенства'':
  \begin{mathpar}
    \inferrule
      {\Gamma \vdash A : \text{type}
        \\ \Gamma \vdash a : A
        \\ \Gamma \vdash b : A}
      {\Gamma \vdash Id_A(a, b) : \text{type}}
    \and
    \inferrule
      {\Gamma \vdash a : A}
      {\Gamma \vdash \text{refl}_A \; a : Id_A(a, a)}
    \and
    \inferrule
      {\Gamma, c : A \vdash t :
        C[a := c, b := c, p := \text{refl}_A \; c]}
      {\Gamma, a : A, b : A, p : Id_A(a, b) \vdash J(t, a, b, p) : C}
  \end{mathpar}
  У него есть свои проблемы, но он наводит на некоторые идеи.
\end{frame}

\section{Системное программирование}

\frame{\sectionpage}

\begin{frame}{Боли безопасного системного программирования}
  \begin{itemize}
    \item В C++ можно делать что угодно, но сама система типов небезопасна;
    \item В Rust безопасная система типов, но порой слишком ограничивающая: иной
      раз ты \textit{можешь доказать}, что правильно обращаешься с памятью, но
      средств языка для этого не хватает.
    \item ATS: введём специальный пропозициональный тип $\text{Access}(T, l)$,
      населённый доказательствами того, что фрагмент кода владеет доступом к
      данным типа $T$ по адресу $l$.
  \end{itemize}
\end{frame}

\begin{frame}{Линейные типы}
  \begin{block}{Линейный Access}
    \begin{itemize}
      \item Доступ нельзя взять и забыть;
      \item Нельзя дублировать доступ.
    \end{itemize}
    С другой стороны, сами ссылки на элемент типа $T$ по адресу $l$ могут
    копироваться беспрепятственно; Access только накладывает ограничения на
    чтение и запись.
  \end{block}
  \begin{block}{Открытая проблема}
    Пока не придумали линейную зависимую систему типов. В данном
    случае можно сгладить углы и зависеть только от типов, на которых не
    наложены субструктурные ограничения.
  \end{block}
\end{frame}

\begin{frame}{Foreign function interface}
  \begin{block}{Задача}
    Делать системные вызовы, а также вызывать код на C через FFI, не теряя
    гарантий системы типов.
  \end{block}
  \begin{block}{Решение}
    Зависимые типы позволяют полностью описать спецификацию системных вызовов (и
    других функций) в типе. Если внешняя функция спецификации не соответствует,
    возможно, получится задействовать механизм ``определения вины'' а-ля gradual
    typing.
  \end{block}
\end{frame}

\begin{frame}{Эффекты и модели стоимости}
  В областях, связанных с системным программированием, программистов часто
  интересует асимптотическая сложность алгоритмов.

  Предлагаемый пайплайн следующий:
  \begin{enumerate}
    \item Сформулировать cost model (время работы? память? внешняя память? кэш?)
    \item Описать DSL для программ, сложность которых мы оцениваем (возможно, в
      си-подобном либо императивном стиле);
    \item Доказывать факты об асимптотической сложности.
  \end{enumerate}

  \textit{Либо: Cost-aware type theory}
\end{frame}

\begin{frame}{Список требований к языку}
  \begin{itemize}
    \item Зависимые типы;
    \item Линейный тип Access;
    \item Производительность;
    \item Скорее всего, отсутствие сборки мусора;
    \item Совместимый с системой типов FFI с C;
    \item Удобство создания императивных DSL;
    \item Средства для оценки асимптотической сложности в стандартной
      библиотеке.
  \end{itemize}
\end{frame}

\end{document}
